%% maestro2tex -- generated table; do not edit by hand.
\begin{table}[htbp]
  \centering
  \caption[{Process, temperature and supply points simulated for the 12-bit R-2R DAC and its supply block: point 1 is the\,\dots}]{Process, temperature and supply points simulated for the 12-bit R-2R DAC and its supply block: point 1 is the typical statistical process point evaluated at nominal, points 2--5 the imported foundry skew corners. Temperature and solver tolerance are not uniform across this run: the two DAC benches sit at \SI{25}{\celsius} at all five points and run \texttt{reltol}~=~\num{1e-6}, \texttt{vabstol}~=~\num{1e-9}, \texttt{iabstol}~=~\num{1e-15}, while the four supply-block benches (load regulation, line regulation, output impedance, startup) sit at \SI{27}{\celsius} at point 1 and \SI{25}{\celsius} at points 2--5 and run the ADE defaults \num{1e-3}, \num{1e-6}, \num{1e-12}; the column below reports the DAC benches. This is a corner run and carries no device mismatch; R-2R linearity is a matching property, so every nonlinearity figure in this section is a systematic term and none of it is a yield statement. These points move the sixteen transistor model rows only (the resistor and capacitor rows stay at their typical sections), so every corner spread in this section is a lower bound (pinned-passive caveat, page~\pageref{p:pinned-passives}).}
  \label{tab:DacR2R12-corners}
  \begin{tabular}{lrrr}
    \toprule
    Point & Process & Temperature & Supply \\
    \midrule
    1 & tt & 25\,\si{\celsius} & 2.5\,\si{\volt} \\
    2 & ff & 25\,\si{\celsius} & 2.5\,\si{\volt} \\
    3 & fs & 25\,\si{\celsius} & 2.5\,\si{\volt} \\
    4 & sf & 25\,\si{\celsius} & 2.5\,\si{\volt} \\
    5 & ss & 25\,\si{\celsius} & 2.5\,\si{\volt} \\
    \bottomrule
  \end{tabular}
\end{table}
